\JOUChapterWithEnglish{浮选柱实验研究}{Experiment Research of Column Flotation}
\label{chap:experiment}

\JOUSectionWithEnglish{浮选柱研究现状}{Present Research of Column Flotation}

快浮装置与旋流-静态微泡浮选柱在煤泥分选领域具有较强的工程适应性。已有工业运行结果表明，合理控制循环矿浆量、入料粒级和背压条件，能够同时改善回收率与精矿灰分指标，并为后续工艺放大提供稳定的数据基础。

\vspace{0.4\baselineskip}

针对不同煤样和矿样开展的系列实验还说明，浮选柱的优势主要体现在细粒级颗粒的界面富集与分选稳定性上。随着在线检测和自动控制手段的引入，实验研究已经逐渐从单因素考察转向多参数耦合分析。

\begin{table}[H]
    \centering
    \JOUTableBicaption{快浮装置工业运行状况}{Industrial operation of fast flotation equipment}
    \label{tab:industrial-operation}
    \renewcommand{\arraystretch}{1.12}
    \begin{tabular}{C{2.4cm}C{2.4cm}C{2.4cm}C{3.2cm}}
        \toprule
        运行工况 & 入料粒级/mm & 精矿灰分/\% & 回收率/\% \\
        \midrule
        工况 A & 0.50--0.074 & 8.42 & 83.50 \\
        工况 B & 0.25--0.045 & 7.96 & 86.12 \\
        工况 C & 0.125--0.045 & 7.61 & 88.03 \\
        \bottomrule
    \end{tabular}
\end{table}

\vspace{0.6\baselineskip}

\begin{table}[H]
    \centering
    \JOUTableBicaption{快浮装置主要技术参数}{Main technical parameter of fast flotation equipment}
    \label{tab:key-parameters}
    \renewcommand{\arraystretch}{1.12}
    \begin{tabular}{C{3.0cm}C{3.2cm}C{4.2cm}}
        \toprule
        参数名称 & 设计值 & 说明 \\
        \midrule
        循环矿浆量 & 6.5\,m\textsuperscript{3}/h & 稳态运行样机 \\
        进气量 & 1.8\,m\textsuperscript{3}/h & 微泡发生器额定值 \\
        柱体背压 & 0.37--1.50\,m 水柱 & 对应三组对比工况 \\
        \bottomrule
    \end{tabular}
\end{table}
